\begin{papercard}{moniz2026}
\cf{Stages} \stagetag{H} (harms and risks as perceived by users of intimate health technologies), \stagetag{S} (definition of ``intimate health technologies'' and ``intimate privacy''), \stagetag{R} (risk-communication agenda). A qualitative HCI study (Int.\ J.\ Human--Computer Interaction 42(15), 2026) using a \emph{design probe} kit with eight Portuguese participants, plus a follow-up interview, analysed by inductive thematic analysis. It is the nearest study in the corpus to the thesis' application domain (cycle/fertility/mental-wellbeing apps) and is mostly a source of user-perceived harms and consent expectations and of a non-traditional method.

\cf{Aim and research questions} To study individual perspectives, experiences and preferences for privacy and consent in digital intimate health technologies, and the concerns and perceived risks attached to them \pc{moniz2026}{2}{This research introduces an empirical study of user perceptions and active privacy management within sensitive technologies}. ``Intimate health technologies'' are tools, apps and devices for personal health in private contexts (sexual and reproductive health, mental well-being, self-care), involving intimate data (sexual orientation, gender, behaviour, reproductive health, medical records) \pc{moniz2026}{2}{Intimate health technologies encompass tools, apps, and devices designed to support personal health in private contexts}. RQ1 how individuals understand, navigate and manage privacy; RQ2 how they decide to consent and what they expect; RQ3 which design features they suggest to improve privacy and consent \pc{moniz2026}{3}{RQ1: How do individuals understand, navigate, and manage privacy in Intimate Health Technologies?}.

\cf{Background claims drawn from prior work} Reviews of women's mHealth apps (Alfawzan et al.\ 2022: 23 apps) found discrepancies in privacy practices, sensitive data collected without adequate protection and poor consent \pc{moniz2026}{4}{Alfawzan et al. (2022) performed a comprehensive scoping review which assessed the privacy, data sharing, and security policies of 23 popular women’s mHealth apps}; menstruation and fertility apps share data without consent and fall short of GDPR (Shipp \& Blasco; Mehrnezhad \& Almeida); BLE vulnerabilities in devices (Cook et al.); recurring consent dark patterns: only an ``Accept'' button, Accept highlighted over Reject, pre-ticked boxes, bundling terms with the policy \pc{moniz2026}{4}{a number of applications adhere to poorly implemented consent practices, such as only presenting the user with one “Accept” option, highlighting Accept over Reject, with pre-ticked boxes, or bundling other terms and conditions together with the privacy policy}.

\cf{Method: design probe} A probe is a packaged set of artefacts and instructions that participants complete over time, designed to gather inspirational insight about values and daily life and to let participants ``find some privacy in the ambiguity of responses'' \pc{moniz2026}{5}{were first introduced by Gaver et al. (1999) as a design-oriented research method}. Kit (a pink envelope of cardboard, paper, markers) with six activities \pc{moniz2026}{6}{The kit consisted of a light pink envelope that contained six activities}: (1) \emph{Ice Breaker} defining ``privacy'' and ``intimate privacy''; (2) \emph{Intimate Diary} (up to seven days of entries); (3) \emph{Intimate Privacy Manifesto}; (4) \emph{Intimate Technology Timeline} of technologies used and feelings about them; (5) \emph{Consent Design}: four questions on defining consent, what influences giving or withholding it, expectations and a redesign sketch; (6) \emph{Futuristic Scenarios} (online): ``Health Technology Revolution: Intimate Data Collection'' (2050, implants, smart clothing) and ``A New Currency: Personal Health Data'' (data traded for benefits) \pc{moniz2026}{6}{The first futuristic scenario was titled “Health Technology Revolution: Intimate Data Collection,”}. Participants kept the kit about 10 days \pc{moniz2026}{6}{Our design probe was an individual set and participants were expected to complete the kit over the course of 10 days}, then a semi-structured 20--25 min interview about impressions, activities and views \pc{moniz2026}{6}{they were invited to participate in a follow-up semi-structured interview, with an estimated duration of 20–25 min}.

\cf{Sample, ethics and analysis} Purposeful sampling via a screening survey in fertility and reproductive-health Discord and Facebook groups in Portugal; women, cisgender, 24--55 years; Lisbon area with in-person kit delivery; no payment; eight participants completed the study \pc{moniz2026}{8}{Eight participants consisting of cisgender women between the ages of 24–55 were recruited and completed the study}. Ethical approval from the IST Lisbon ethics committee, printed informed-consent forms \pc{moniz2026}{8}{Ethical approval was obtained from the University of Lisbon, Instituto Superior Tecnico’s Ethics Committee prior to the commencement of the study}. Data: handwritten responses, digital submissions, interview audio, transcribed, translated to English and checked; \emph{inductive thematic analysis} in Braun \& Clarke's sense, coding iteratively, grouping and ``letting the themes emerge'' \pc{moniz2026}{8}{We conducted a thematic analysis and applied an inductive approach by coding data iteratively, grouping it, and letting the themes emerge}. Word counts per activity are tabulated (Table 2); the diary drew the most text. The authors list small sample size, self-selection of privacy-interested participants and the effort/non-payment as limits \pc{moniz2026}{8}{The study used a small sample size of eight participants, which may not reflect the diverse target population}.

\cf{Findings (themes)} \emph{Understanding of privacy}: all eight defined privacy as a ``right'', as control over disclosure and as a personal safe space; \emph{intimate privacy} was distinguished as concerning body, sexuality and the most personal aspects of life, deserving higher care \pc{moniz2026}{10}{Differentiating from privacy, intimate privacy was considered more specifically related to their body, sexuality, and the most personal aspects of their lives}. \emph{``The rush'', transparency and consent}: consent decisions are driven by hurry and convenience (``click agree and move on''), by perceived trustworthiness and reviews; participants want transparency about ``what, who, why, when, and where'' data will be used, plain language, and consent terms that ``reflect what happens in reality'' \pc{moniz2026}{11}{P2 talked about “the rush” as the main influencing factor when providing consent on digital platforms}; they proposed personalised, interactive or multimedia consent, interactive pop-ups, chatbots, granular options rather than ``only one main `I consent' box'', periodic notifications and easy revocation \pc{moniz2026}{12}{consent mechanisms should offer options to choose from and select based on your preferences, not only containing one main “I consent” box}. \emph{Burdens and invisible threats}: in the futuristic scenarios participants weighted risks above benefits \pc{moniz2026}{13}{Through the futuristic scenarios, they more prominently highlighted the risks and considered them to be more important than the benefits}; emotional strain of constant tracking, anxiety and a feeling of failure when dates were missed, anxiety from pregnancy-related suggestions after a delayed period; a participant on Flo noticed from 2023 menstruation-related and mood-related content on Instagram Reels and TikTok after entering such data in the app and called it a privacy violation she nonetheless accepts \pc{moniz2026}{13}{participants expressed feelings of anxiety (P2, Intimate Diary) and even failure (P6, Intimate Diary) because they had missed dates in their tracking}. Commodification: for the data-as-currency scenario most rated comfort 1 or 2 of a scale (Table 4: S.2 mostly ``very uncomfortable''), arguing that ``health is not tradable'', that it exploits vulnerable people and may worsen inequality; concern about inescapable surveillance, discrimination; more trust in conventional medical systems and professionals \pc{moniz2026}{14}{Participants highlighted having higher trust in conventional medical systems and professionals}. \emph{Interviews}: favourite activities Manifesto, Futuristic, Diary, Timeline, Consent, Ice-Breaker; 7/8 said the kit changed their thinking, two planned to delete unused period-tracking apps \pc{moniz2026}{15}{the most common answers, in ranked order, were: 1. Privacy Manifesto, 2. Futuristic Activity, 3. Intimate Diary, 4. Timeline, 5. Consent, and 6. Ice-Breaker}.

\cf{Discussion: design considerations and a risk-communication agenda} Three considerations built on participants' ideas and cited studies: (1) short videos (Obar et al.: brief live-action under a minute preferred) \pc{moniz2026}{16}{participants favored brief (less than one minute) live-action videos over longer, animated ones}; (2) chatbots (Rumi for informed consent); (3) interactive tools and granular privacy options (Reinhardt et al.'s VIPP; Poli-see), noting that most intimate-health apps use a ``take-it-or-leave-it'' approach \pc{moniz2026}{17}{most intimate health applications also being part of the “take-it-or-leave-it” approach}. Since any new format can lead back to habituation, they argue for \emph{risk communication}: informing users of the specific risks of processing in an understandable, actionable way \pc{moniz2026}{18}{Risk communication involves informing users about the specific risks associated with data processing activities in a way that is easily understandable and actionable}. Five principles: contextual sensitivity (notices tailored to the data use and to user groups, e.g.\ by country of birth, religion or potential for harm) \pc{moniz2026}{18}{Privacy notices and consent forms should be tailored to the specific context of data use, highlighting risks relevant to different user groups}, visual and interactive elements (e.g.\ personal data-flow map), ongoing communication (periodic updates, one-page consent review), personalised risk information (dashboard with a dynamic risk score, ``privacy scenarios'') \pc{moniz2026}{19}{Offer personalized risk assessments based on users’ specific data and context}, and clear simple language.

\cf{Limits (card)} Eight women from one country, self-selected and privacy-interested, so no prevalence or generalisation; no reported coder count, codebook or agreement for the thematic analysis; findings are self-reports and speculative-scenario reactions rather than observed behaviour; proposals (chatbot, risk score, video) are untested in this study; the claim of ``global relevance'' rests on shared apps, not on data; the paper does not analyse actual app behaviour; no distinction between app types (cycle trackers vs.\ mental-health apps).

\cf{Use for this thesis} (1) Domain-matched evidence of perceived harms (emotional strain from tracking, targeted content after entering intimate data, commodification, discrimination), useful as candidate nodes and as user-voice evidence in the harm tree (stage H) but not as prevalence. (2) A source for the construct ``intimate privacy'' for stage S definitions and for the claim that consent in these apps is dominated by haste and take-it-or-leave-it designs, to be checked against what the walkthrough observes. (3) The list of consent dark patterns (single Accept, Accept highlighted, pre-ticked boxes, bundling) is a ready checklist for the walkthrough codebook, citing Moniz et al.\ as a secondary source and the primary regulatory/empirical literature for definitions. (4) Research-ethics template for sensitive topics (participants control what they share, consent forms, withdrawal). (5) The risk-communication agenda is a candidate output type for the thesis' recommendations, but must be framed as a proposal and not as validated.
\end{papercard}
